
\documentclass[]{beamer}
\usetheme{Madrid}
\usecolortheme{default}

\usepackage[utf8]{inputenc}
\usepackage{amsmath, amssymb}
\usepackage{hyperref}

\title[Shallow-water model]{Hands-on shallow-water modelling with Python}
\author{}
\date{\today}

\begin{document}

%------------------------------------------------
\begin{frame}
  \titlepage
\end{frame}
%------------------------------------------------



\begin{frame}{What does this shallow-water model do?}
  \begin{itemize}
    \item 2D shallow-water equations on an \textbf{Arakawa C-grid}
    \item Rectangular basin on an \textbf{f-plane or $\beta$-plane}
    \item Options for:
    \begin{itemize}
      \item Linear / nonlinear momentum equations
      \item Rayleigh friction and viscosity
      \item Range of forcings (wind, tides, storms, etc.)
    \end{itemize}
    \item Designed to be:
    \begin{itemize}
      \item \textbf{Small and readable}: good for learning and experimenting
      \item \textbf{Modular}: separate grid, parameters, forcings, numerics, plots
    \end{itemize}
    \item Comes with:
    \begin{itemize}
      \item A \textbf{public Python API} for scripts and notebooks
      \item Example \textbf{Jupyter notebooks} for classical ocean-dynamics problems
    \end{itemize}
  \end{itemize}
\end{frame}



\begin{frame}[shrink=8]{Tour of the repository}
  \begin{block}{Top level}
    \begin{itemize}
      \item \texttt{README.md} -- quick start, installation, examples
      \item \texttt{notebooks/} -- study cases (wind gyre, Rossby waves, gravity waves, tides, storms\dots)
      \item \texttt{pyproject.toml}, \texttt{setup.cfg} -- packaging info
    \end{itemize}
  \end{block}
  \begin{block}{Main code (Python package)}
    \texttt{src/shallowwater/}
    \begin{itemize}
      \item \texttt{\_\_init\_\_.py} -- exposes the public API
      \item \texttt{grid.py} -- grid creation, metric factors (C-grid)
      \item \texttt{params.py} -- \texttt{ModelParams} and physical/numerical parameters
      \item \texttt{initial.py} -- initial conditions (Gaussian bumps, balanced states, \dots)
      \item \texttt{forcing.py} -- wind stress, tidal forcing, storms, \dots
      \item \texttt{tendencies.py}, \texttt{time\_stepping.py} -- core numerical scheme
      \item \texttt{visualize.py} -- helper functions for plots and animations
    \end{itemize}
  \end{block}
  \begin{block}{Mental map}
    \textbf{User-facing}: \texttt{\_\_init\_\_.py}, \texttt{notebooks/}, \texttt{visualize.py}\\
    \textbf{Under the hood}: \texttt{grid.py}, \texttt{tendencies.py}, \dots
  \end{block}
\end{frame}



\begin{frame}[shrink=8]{Public API cheat sheet}
  \begin{block}{Model configuration}
    \begin{itemize}
      \item \texttt{ModelParams(...)}\\
        Physical: \texttt{H}, \texttt{g}, \texttt{rho}, \texttt{f0}, \texttt{beta}, \texttt{y0}\\
        Damping/viscosity: \texttt{r}, \texttt{Ah}, \dots\\
        Switches: \texttt{linear=True/False}, toggle advection, \dots
      \item \texttt{make\_grid(Nx, Ny, Lx, Ly)} builds the computational grid
    \item[]
    \end{itemize}
  \end{block}
  \begin{block}{Initial conditions and forcing}
    \begin{itemize}
      \item \texttt{setup\_initial\_state(grid, params, mode=..., ...)}\\
        e.g. \texttt{mode="gaussian\_bump"}, \texttt{mode="balanced\_jet"}
      \item Forcing functions like \texttt{zero\_forcing}, \texttt{wind\_gyre\_forcing}, \dots
      \item You can define your \textbf{own forcing function} with the same signature
    \end{itemize}
  \end{block}
  \begin{block}{Time stepping and visualisation}
    \begin{itemize}
      \item \texttt{compute\_dt\_cfl(grid, params, cfl=0.5)} suggests a stable time step
      \item \texttt{run\_model(tmax, dt, grid, params, forcing\_fn, ic\_fn, ...)} integrates the equations
      \item Visualisation helpers: \texttt{animate\_eta(...)}, \texttt{plot\_forcings(...)}, Hovmöller/spectral plots
    \end{itemize}
  \end{block}
\end{frame}




\begin{frame}[fragile,shrink=8]{Where to find the \texttt{shallowwater} package}
\begin{itemize}
  \item The \texttt{shallowwater} code is hosted on GitHub:
\begin{verbatim}
https://github.com/fabien-roquet/shallowwater.git
\end{verbatim}

  \item We will use the tagged version:
\begin{verbatim}
tag: v0.1.1
\end{verbatim}

  \item \textbf{What is GitHub?}
    \begin{itemize}
      \item Online platform to store and share code.
      \item Uses \texttt{git} for version control (history of changes).
      \item Makes it easy to collaborate (issues, pull requests).
    \end{itemize}

  \item The package is also on the Python Package Index (PyPI) and can be
        installed with:
\begin{verbatim}
pip install shallowwater
\end{verbatim}

  \item In this lab you will:
    \begin{itemize}
      \item Download the GitHub repository locally (to see the code, docs, examples).
      \item Install the library via \texttt{pip install shallowwater} (to use it in Python).
    \end{itemize}
\end{itemize}
\end{frame}



\begin{frame}[fragile,shrink=8]{Downloading the repository with \texttt{git}}
\begin{itemize}
  \item If you have \texttt{git} installed (Linux, macOS, or Windows):

    \begin{itemize}
      \item Open a terminal (or a terminal inside JupyterLab).
      \item Go to a folder where you want to store the code, e.g.:
\begin{verbatim}
cd ~/projects
\end{verbatim}
      \item Clone the repository:
\begin{verbatim}
git clone https://github.com/fabien-roquet/shallowwater.git
cd shallowwater
\end{verbatim}
      \item Switch to the tagged version \texttt{v0.1.1}:
\begin{verbatim}
git checkout v0.1.1
\end{verbatim}
    \end{itemize}

  \item You now have a local copy of:
    \begin{itemize}
      \item Python source code in \texttt{src/shallowwater/}
      \item Example notebooks in \texttt{notebooks/}
      \item Documentation and configuration files
    \end{itemize}
\end{itemize}
\end{frame}



\begin{frame}[fragile,shrink=8]{Downloading the repository without \texttt{git}}
\begin{itemize}
  \item If you do not have \texttt{git} (common on Windows), you can use a ZIP file instead:
    \begin{itemize}
      \item Open the GitHub page in your web browser.
      \item Click \textbf{"Code"} $\rightarrow$ \textbf{"Download ZIP"}.
      \item Save the ZIP file to your computer.
      \item Unzip it to a folder, for example:
\begin{verbatim}
Documents/shallowwater/
\end{verbatim}
    \end{itemize}

  \item After unzipping:
    \begin{itemize}
      \item You can see the same contents as in the git clone
            (code, notebooks, etc.).
      \item You can point JupyterLab’s file browser to that folder to open files.
    \end{itemize}

  \item Optional: check the version information (in the \texttt{README} or
        release notes) to confirm it corresponds to version \texttt{0.1.1}.
\end{itemize}
\end{frame}


\begin{frame}[fragile,shrink=8]{Installing \texttt{shallowwater} from PyPI}
\begin{itemize}
  \item Once Python and \texttt{pip} are available, install the package with:
\begin{verbatim}
pip install shallowwater
\end{verbatim}

  \item This fetches the released version from PyPI and makes it available as a
        normal Python library.

  \item You can quickly verify that the installation works:
\begin{verbatim}
python -c "import shallowwater"
\end{verbatim}

  \item After installation, you can \texttt{import shallowwater} in any
        Python script or notebook, regardless of where the GitHub repo sits.

  \item No special virtual environment is strictly required here:
    \begin{itemize}
      \item Only \texttt{numpy} is needed.
      \item \texttt{numba} (if installed) is used optionally for speed.
    \end{itemize}
\end{itemize}
\end{frame}



\begin{frame}[fragile,shrink=8]{Using the \texttt{lab\_tsunami.ipynb} notebook}
\begin{itemize}
  \item The lab material is provided separately on Canvas as:
\begin{verbatim}
lab_tsunami.ipynb
\end{verbatim}

  \item Steps:
    \begin{itemize}
      \item Download \texttt{lab\_tsunami.ipynb} from Canvas.
      \item Save it in a convenient folder, for example inside your local
            \texttt{shallowwater} directory.
    \end{itemize}

  \item Start JupyterLab (from your usual method) and:
    \begin{itemize}
      \item Use the file browser to navigate to the folder where
            \texttt{lab\_tsunami.ipynb} is stored.
      \item Open the notebook by clicking on it.
    \end{itemize}

  \item The notebook will:
    \begin{itemize}
      \item \texttt{import shallowwater} (installed via \texttt{pip install shallowwater}).
      \item Guide you through running a tsunami-like simulation.
      \item Ask you to experiment with key parameters:
        \begin{itemize}
          \item water depth \texttt{H}
          \item initial Gaussian disturbance
        \end{itemize}
    \end{itemize}
\end{itemize}
\end{frame}



\end{document}
